\section{Background \& Related Work}
\label{sec:related-work}

Table~\ref{tab:benchmark-comparison} summarizes how RealVuln relates to prior
benchmarks across seven desirable properties. We discuss each category of prior
work below, highlighting the specific gaps that motivated our design.

%% ── Comparison table ─────────────────────────────────────────────────────────
\begin{table*}[t]
\caption{Comparison of SAST benchmarks. \emph{Real Code}: uses unmodified
  application source, not synthetic snippets. \emph{Open Data}: ground truth
  and scanner results publicly available. \emph{Open Scoring}: scoring code
  released for reproduction. \emph{FP Testing}: explicitly measures false
  positives. \emph{Multi-Scanner}: evaluates $\geq$3 scanners head-to-head.
  \emph{LLM Scanners}: includes LLM-based tools. \emph{Multi-Lang}: covers
  more than one programming language.}
\label{tab:benchmark-comparison}
\centering
\footnotesize
\begin{tabular}{@{}lccccccc@{}}
\toprule
\textbf{Benchmark} & \textbf{Real} & \textbf{Open} & \textbf{Open} & \textbf{FP} & \textbf{Multi-} & \textbf{LLM} & \textbf{Multi-} \\
 & \textbf{Code} & \textbf{Data} & \textbf{Scoring} & \textbf{Test} & \textbf{Scanner} & \textbf{Scan.} & \textbf{Lang} \\
\midrule
OWASP Benchmark~\cite{owaspbenchmark}
  & $\times$ & \checkmark & \checkmark & \checkmark & \checkmark & $\times$ & $\times$ \\
Juliet Test Suite~\cite{boland2012juliet}
  & $\times$ & \checkmark & $\times$ & $\times$ & $\times$ & $\times$ & $\times$ \\
SastBench~\cite{feiglin2026sastbench}
  & \checkmark & $\times$ & $\times$ & \checkmark & $\times$ & $\times$ & \checkmark \\
ZeroPath/XBOW~\cite{zeropath2024}
  & $\times$ & $\sim$ & $\times$ & \checkmark & \checkmark & \checkmark & $\times$ \\
Fluid Attacks~\cite{fluidattacks2025}
  & \checkmark & $\times$ & $\times$ & \checkmark & \checkmark & $\times$ & $\times$ \\
CASTLE~\cite{dubniczky2025castle}
  & $\times$ & \checkmark & \checkmark & $\times$ & \checkmark & \checkmark & \checkmark \\
Cycode~\cite{cycode2024benchmark}
  & \checkmark & $\times$ & $\times$ & $\times$ & \checkmark & $\times$ & \checkmark \\
DryRun~\cite{dryrun2025sast}
  & \checkmark & $\sim$ & $\times$ & $\times$ & \checkmark & $\times$ & \checkmark \\
\midrule
\textbf{RealVuln (ours)}
  & \checkmark & \checkmark & \checkmark & \checkmark & \checkmark & \checkmark & $\times$ \\
\bottomrule
\end{tabular}
\end{table*}

%% ── §1  Synthetic benchmarks ────────────────────────────────────────────────
\paragraph{Synthetic benchmarks.}
The OWASP Benchmark~\cite{owaspbenchmark} (${\sim}$2{,}740 synthetic Java servlets) and the NIST Juliet Test Suite~\cite{boland2012juliet} (single-CWE C/C++/Java programs) are the most cited SAST evaluation suites. Both use short, self-contained test cases that let pattern-matching tools score well without understanding cross-file data flow or framework semantics, limiting their relevance to whole-application analysis.

%% ── §2  Unreleased datasets ─────────────────────────────────────────────────
\paragraph{Unreleased datasets.}
SastBench~\cite{feiglin2026sastbench} assembles 2{,}737 ground-truth samples across 38 languages from real CVEs, scored with the Matthews Correlation Coefficient. Its dataset is not publicly released, and it evaluates agentic \emph{triage} of SAST alerts rather than raw scanner detection accuracy. SastBench is complementary to RealVuln. It measures post-scan triage that we deliberately exclude, but cannot serve as an open benchmark for detection performance.

%% ── §3  Vendor benchmarks ───────────────────────────────────────────────────
\paragraph{Vendor benchmarks.}
ZeroPath~\cite{zeropath2024} published benchmark code and SARIF results as a public GitHub fork of the XBOW validation benchmarks, testing against major SAST vendors, a genuine step toward transparency, though without reusable scoring infrastructure or ground-truth labels. Fluid Attacks~\cite{fluidattacks2025} tested 36 scanners against a large JS/TS application (the best achieved 22.7\% recall), but neither the source nor the ground truth is public. Cycode~\cite{cycode2024benchmark} ran Bearer CLI against open-source projects and published raw classification data, but the methodology is not fully documented for reproduction. DryRun~\cite{dryrun2025sast} tested against public intentionally-vulnerable applications with per-vulnerability breakdowns, providing partial reproducibility. The shared limitation across all vendor benchmarks is a structural conflict of interest: each vendor selects the test conditions, interprets results, and lacks a unified framework for independent re-scoring.

%% ── §4  Dataset quality ─────────────────────────────────────────────────────
\paragraph{Dataset quality and label noise.}
Automated vulnerability datasets suffer from widespread label noise. Ding et al.~\cite{ding2024primevul} showed that a 7B model achieves 68\% F1 on BigVul but only 3\% on their cleaned PrimeVul dataset; CleanVul~\cite{li2025cleanvul} found over 40\% of entries in existing datasets mislabeled. CVEFixes~\cite{bhandari2021cvefixes} and ReposVul~\cite{wang2024reposvul} inherit the same tangled-commit problems. RealVuln takes the opposite position: 796 hand-labeled findings across 26 repositories, accepting lower scale in exchange for label quality that can be audited line by line.

%% ── §5  Agentic and LLM benchmarks ─────────────────────────────────────────
\paragraph{Agentic and LLM-oriented benchmarks.}
SWE-Bench~\cite{jimenez2024swebench} introduced the approach of evaluating AI agents on real GitHub issues with reproducible harnesses. SWE-Bench Verified~\cite{openai2025swebenchverified} further demonstrated the value of human-validated labels. CVE-Bench~\cite{zhu2025cvebench} measures whether agents can exploit real CVEs in containerized apps, highlighting knowledge-cutoff confounds. CASTLE~\cite{dubniczky2025castle} evaluates 25 tools (13 static analyzers, 10 LLMs, and 2 formal verification tools) on CWE detection but uses synthetic micro-benchmarks. Yildiz et al.~\cite{yildiz2025jitvul} found that agentic approaches outperformed function-level prompting for vulnerability detection but noted that agents still exhibit inconsistencies requiring further refinement.

%% ── §6  Honest self-critique ────────────────────────────────────────────────
\paragraph{Positioning and limitations of RealVuln.}
No prior benchmark systematically compares all three scanner
categories (Rule-Based SAST, General-Purpose LLMs, and Security-Specialized
systems) on identical ground truth. As Table~\ref{tab:benchmark-comparison}
shows, RealVuln is the only benchmark that simultaneously uses real code,
releases all data and scoring code, tests for false positives, and includes
scanners from all three categories.
We are forthright about its limitations in v1. First, RealVuln covers
Python only. This enables depth (Django, Flask, FastAPI, and raw Python
idioms are all represented), it cannot speak to scanner performance on memory-unsafe
languages or compiled ecosystems. Second, our ground truth covers only
OWASP-style application vulnerabilities (injection, XSS, authentication
flaws), but we do not yet include supply-chain, configuration, or
infrastructure-as-code findings. Third, we are ourselves a vendor (Kolega.Dev),
which creates the same conflict of interest we criticize above. We mitigate
this by open-sourcing \emph{everything}: ground truth labels, all scanner
result files, the scoring pipeline, and the interactive dashboard, so that any
researcher can audit, challenge, or extend our results. We view transparency
not as a defense against bias but as an invitation for the community to find it.
